\documentclass[10pt,a4paper]{article}

\usepackage[T1]{fontenc}
\usepackage[utf8]{inputenc}
\usepackage{lmodern}
\usepackage[a4paper,margin=2cm]{geometry}
\usepackage{array,booktabs,longtable}
\usepackage{amsmath,amssymb}
\usepackage[dvipsnames]{xcolor}
\usepackage[round,authoryear]{natbib}
\usepackage[colorlinks=true,linkcolor=MidnightBlue,citecolor=MidnightBlue,urlcolor=MidnightBlue]{hyperref}
\usepackage{doi}
\usepackage{enumitem}

% Fit labels at the start of each description
\newcommand{\fitlabel}[2]{\textcolor{#1}{\textbf{#2}}}
\newcommand{\confirms}{\fitlabel{ForestGreen}{Confirms:}}
\newcommand{\consistent}{\fitlabel{RoyalBlue}{Consistent:}}
\newcommand{\method}{\fitlabel{Plum}{Method:}}
\newcommand{\own}{\fitlabel{BurntOrange}{Own finding:}}
\newcommand{\abs}{$^\dagger$}   % content only checked via abstract or summary

\newcolumntype{L}[1]{>{\raggedright\arraybackslash}p{#1}}
\setlength{\tabcolsep}{4pt}
\renewcommand{\arraystretch}{1.25}

% One table per topic: our insight | supporting paper | short description
\newenvironment{insighttable}{%
  \small
  \begin{longtable}{L{0.36\textwidth} L{0.17\textwidth} L{0.41\textwidth}}
  \toprule
  \textbf{Our insight} & \textbf{Supporting paper} & \textbf{Short description} \\
  \midrule
  \endhead
  \bottomrule
  \endfoot
}{%
  \end{longtable}
}

\title{\textbf{MODIS $\rightarrow$ VIIRS harmonisation for Iberia}\\[2pt]
       \large Our insights and the research behind them}
\author{La-Brasa team, NASA Space Apps Challenge 2026}
\date{15 November 2026}

\begin{document}
\maketitle

\noindent
This overview lists what we found while harmonising the NASA FIRMS hotspots of MODIS (Terra, Aqua) and VIIRS
(SNPP, NOAA-20, NOAA-21) for the Iberian Peninsula, and which research papers support each finding.
We found the insights in our own data (\texttt{harmonisation/data\_analysis/}). The papers show the same effect
elsewhere or justify the method we use. Each description starts with how well the paper fits:

\begin{itemize}[nosep,leftmargin=1.5em]
  \item \confirms\ the paper reports the same effect, often with a similar number.
  \item \consistent\ the paper points the same way but measures something slightly different.
  \item \method\ the paper is the standard source for a method we use.
  \item \own\ no paper found. The evidence is our own validation: 42 KPIs in
        \texttt{data/harmonised/} \texttt{kpi\_report.json}.
\end{itemize}

\medskip\noindent
All numbers refer to Iberia, vegetation fires only (\texttt{type} $= 0$), calibration window 02/2012--12/2021.
Bibliographic data of all papers was checked against Crossref. Papers marked \abs\ were only read as an abstract
or summary; read the full text before relying on details.

% ---------------------------------------------------------------------------
\section{Sensors: why MODIS and VIIRS differ}

\begin{insighttable}
VIIRS SNPP reports $6.0\times$ as many detections as Aqua and $\approx 3.0\times$ as many as Aqua + Terra
(by day $4.1\times$, at night $11\times$).
& \citet{vadrevu2018}
& \confirms\ VIIRS detected $6.5\times$ as many fires as Aqua MODIS in an agricultural region of India. \\

VIIRS sees many more small fires because of its 375\,m pixels.
& \citet{schroeder2014}
& \confirms\ describes the VIIRS 375\,m algorithm; it detects far more fire pixels than the coarser 750\,m
product, especially at night. \\

The difference is \emph{which} fires are detected: Aqua detects a VIIRS fire cluster with 50\,\% probability
only from $\approx 34$\,MW (day) / 21\,MW (night). Aqua sees 24\,\% of the clusters, but they contain 83\,\% of the FRP.
& \citet{giglio2016}; \citet{schroeder2008}; \citet{sofiev2026}\abs
& \consistent\ MODIS detection probability depends mainly on fire size; detection limit 3.7\,MW (night, nadir)
up to 40\,MW (day, off-nadir). Our 50\,\% point in MW is our own fit. \\

The daily FRP sum is comparable between sensors (SNPP/Aqua $\approx 1.04$); the factor does not transfer
worldwide (0.9--3.3).
& \citet{li2018}; \citet{li2020}\abs
& \confirms\ FRP of fire clusters is very comparable; differences $\approx 20$\,\% on a $1^\circ$ grid that vary
by region. In Africa MODIS daily FRP is at least 42.8\,\% too low. \\

Single-pixel FRP is not comparable (a VIIRS pixel is $0.45\times$ as strong as a MODIS pixel of the same fire).
Only sums are.
& \citet{li2018}; \citet{zheng2026}\abs
& \consistent\ per-pixel FRP differs between sensors and MODIS FRP depends strongly on the view angle;
Zheng et al.\ build a view-angle correction for MODIS FRP. \\

\texttt{brightness} is not comparable: the VIIRS channel saturates at 367\,K (14\,\% of day detections sit
exactly on this value). Dropped.
& \citet{schroeder2014}
& \confirms\ names saturation of the middle-infrared channel as a known limitation of VIIRS. \\

VIIRS confidence ``h'' only means ``saturated'' (all h rows are at 367\,K); MODIS confidence rises with
brightness. The classes are not comparable.
& --
& \own\ fits the saturation note in \citet{schroeder2014}, but no paper says it directly. \\

Overpass time matters: fires burn strongest in the afternoon, so Aqua/SNPP ($\approx$ 13:30) and Terra
($\approx$ 10:30) see different amounts. FRP before 2012 therefore comes from Aqua alone (H1).
& \citet{giglio2006}; \citet{freeborn2011}; \citet{andela2015}
& \confirms\ strong diurnal fire cycle; polar orbiters see only snapshots of it, and the overpass time
changes what is measured. \\

Orbit drift: in 2026 Aqua flies over at $\approx$ 14:53 instead of 13:05 local time (+108\,min), Terra at 9:21
instead of 11:01 ($-101$\,min); SNPP is stable. Therefore no MODIS after 2021.
& \citet{giglio2025}
& \consistent\ confirms the dates (last Terra manoeuvre 2022, Aqua 2021) and the effect on diurnal sampling;
the impact on fire products is ``under active investigation''. \textbf{Our drift in minutes is our own number.} \\

Two identical VIIRS instruments (NOAA-20 vs.\ SNPP) differ by $\times 1.18$--$1.27$ per $1^\circ$ cell and month.
The mean of all active VIIRS satellites hits a held-out satellite better than any single one (N3).
& \citet{csiszar2018}\abs; \citet{zhou2023}\abs
& \consistent\ overpasses $\approx 50$\,min apart, no perfect agreement expected; FRP of SNPP and NOAA-20
consistent overall. The size of the noise floor and the benefit of the mean are our own result. \\
\end{insighttable}

% ---------------------------------------------------------------------------
\section{Data cleaning}

\begin{insighttable}
Static sources (industry, gas flares) are 32\,\% of SNPP night detections but only 4\,\% for Aqua. They distort
the sensor comparison and are removed via \texttt{type}.
& \citet{giglio2016}; \citet{firms_type}
& \method\ \texttt{type} separates vegetation fires from static land sources and offshore sources.
The 32\,\% vs.\ 4\,\% is our own number. \\

In the near-real-time data \texttt{type} is missing. A location mask learned from the archive replaces it
(recall $\approx 93$\,\%, 1.5\,\% of real fires removed).
& \citet{firms_type}; \citet{ma2020}\abs; \citet{gao2026}\abs
& \consistent\ static sources are found through persistence at the same place over years (FIRMS: 400\,m grid,
$\geq 5$ detections per year). Our thresholds and quality numbers are our own. \\
\end{insighttable}

% ---------------------------------------------------------------------------
\section{\texorpdfstring{Calibration (MODIS $\rightarrow$ VIIRS equivalent)}{Calibration (MODIS to VIIRS equivalent)}}

\begin{insighttable}
The factor is the sum of VIIRS divided by the sum of MODIS over the whole window, not the mean of daily ratios.
& \citet{cochran1977}\abs; \citet{gregoire2009}
& \method\ the ratio of sums (ratio estimator) is much less biased than the mean of ratios, and days with
0 MODIS detections cause no division by 0. \\

The factor differs by region (2.5--5.1 against Aqua + Terra) but is stable over the years, because the
typical fire size differs by region.
& \citet{li2018}
& \consistent\ the MODIS--VIIRS relationship varies between regions and spatial scales. \\

Simple beats complex: of 17 conversions tested (factors per month, day/night, Poisson models, gradient boosting,
quantile mapping, FRP-class weights), one factor per region hit the VIIRS class best (64\,\%).
& \citet{hand2006}\abs
& \consistent\ general argument that simple models often do as well as complex ones on new data.
\textbf{The comparison itself is our own.} \\

A sensor change must not show up as a jump in the series (without conversion the mean would seem to jump
$2.2\times$ in 2012).
& \citet{venema2012}
& \method\ homogenisation of climate time series after instrument changes; the established framework for
our problem. \\

Related approach from the literature.
& \citet{zheng2026}\abs
& \consistent\ fuses Aqua MODIS and SNPP VIIRS FRP (2012--2024) with corrections. FRP only, no detection
counts, no regional factors. Closest related work to compare against. \\
\end{insighttable}

% ---------------------------------------------------------------------------
\section{Error model and aggregation}

\begin{insighttable}
The $0.25^\circ$ grid is too fine for comparison: $R^2$ (log) 0.37 per cell and month, 0.73 per NUTS-2 region
and month, 0.89 per year. The error shrinks with area and time step (day $\times 2.08$, month $\times 1.5$,
year $\times 1.12$).
& \citet{openshaw1984}\abs; \citet{li2018}
& \consistent\ results depend on the size of the aggregation unit (modifiable areal unit problem); the
MODIS--VIIRS relationship changes with spatial scale. Our error per scale is our own measurement. \\

The error depends on the size of the number: $\times 3.0$ at 0, $\times 2.3$ at 1--5, $\times 1.14$ above 500.
A fixed factor ($\times 1.46$) would be far too optimistic for small numbers.
& \citet{gehrels1986}\abs
& \method\ counting statistics: random counts scatter by about $\sqrt{n}$, so the relative error is large
for small counts. \\

Before 2012, MODIS reports nothing on $\approx 60$\,\% of region-days on which VIIRS sees fire. A 0 before 2012
often only means ``MODIS saw nothing''.
& \citet{li2020}\abs
& \consistent\ MODIS misses many low-intensity fires. The 60\,\% is our own number. \\

Big fire days (FRP sum $\geq 100$\,MW per region and day): MODIS finds 62--67\,\%; the extreme class
($\geq 500$\,MW) is right in 72--76\,\%. Medium days are poor.
& \citet{sofiev2026}\abs; \citet{giglio2016}
& \consistent\ MODIS detects powerful fires with $\approx 78$\,\% probability worldwide; detection rises with
fire size. \\

Climatology: the converted MODIS seasonal profile matches VIIRS with $R^2 = 0.94$ after smoothing over
$\pm 7$ days, only 0.75 without.
& \citet{giglio2006}; \citet{zhang2005}\abs
& \consistent\ Terra and Aqua agree well on peak month and season length; a window around each calendar day
is the standard way to stabilise day-of-year statistics. The $\pm 7$ days is our own choice from the test. \\
\end{insighttable}

% ---------------------------------------------------------------------------
\section{Evaluation methods}

\begin{insighttable}
Leave-one-year-out cross-validation; error bands fitted on one half of the years and checked on the other.
& \citet{roberts2017}
& \method\ with time-dependent data, random cross-validation is too optimistic; blocking by time (here: by year)
is the right choice. \\

Error given as a factor, $\exp\left|\log(\text{prediction}/\text{measurement})\right|$ (e.g.\ D1 ``$\times 2.08$'').
& \citet{morley2018}
& \method\ log accuracy ratio and median symmetric accuracy, a symmetric error measure for data that spans
orders of magnitude. \\

Bands must hold their promise: the 50\,\% band contains the measurement in $\approx 50$\,\% of cases, the 90\,\% band
in $\approx 90$\,\% (U1, U2).
& \citet{gneiting2007}; \citet{lei2018}\abs
& \method\ calibration of probabilistic predictions; empirical quantiles on held-out data correspond to split
conformal prediction. \\

Agreement of classes (month class, big fire day) measured with Cohen's kappa.
& \citet{cohen1960}\abs; \citet{pontius2011}
& \method\ (with caveat): kappa is standard but contested in remote sensing; Pontius \& Millones recommend
also reporting quantity and allocation disagreement. \\
\end{insighttable}

\noindent
General overview of active fire remote sensing: \citet{wooster2021}.

\bibliographystyle{plainnat}
\bibliography{references}

\end{document}
